\documentclass[12pt]{article}
\usepackage[utf8]{inputenc}
\usepackage[T1]{fontenc}
\usepackage{hyperref}
\title{Tupar\'i Language Resources for K-12 Curriculum Development}
\author{Compiled by Claude Sonnet 5 (claude-sonnet-5) for Andr\'as Kornai}
\date{}
\begin{document}
\maketitle

\section{Introduction}
Tupar\'i (ISO 639-3: tpr; Glottocode tupa1250) is a Tupian language of Rond\^onia, Brazil, with roughly 340--350 speakers among about 600 ethnic members. It is endangered. Search found a good academic grammar and a corpus, but no school materials.

\section{Primary Educational Resources}
No primers or textbooks were found in search results.

\section{Linguistic and Reference Materials}
\begin{itemize}
\item Adam Singerman (2018), dissertation on the morphosyntax of Tupar\'i, University of Chicago, the most extensive grammar available: \url{https://knowledge.uchicago.edu/records/s4z1e-xtf32}
\item Summary of Singerman (2018): \url{https://websites.umass.edu/phonolist/?p=2559}
\item Poliana Alves produced a bilingual dictionary (2004) and a phonology (1991), as reported on Wikipedia. I did not find a link to the dictionary itself: \url{https://en.wikipedia.org/wiki/Tupari_language}
\item Comparative paper on non-verbal predicates in the Tupar\'i languages (Museu Goeldi repository): \url{https://repositorio.museu-goeldi.br/jspui/bitstream/mgoeldi/2709/1/Predicados%20nominais%20.....pdf}
\end{itemize}

\section{Religious and Community Resources}
None found.

\section{Digital and Audio Resources}
\begin{itemize}
\item The Tupar\'i corpus by Adam Singerman, hosted at AILLA (UT Austin): \url{https://islandora-ailla.lib.utexas.edu/islandora/object/ailla%3A269477}
\end{itemize}

\section{Curriculum Development Context}
The grammar, the dictionary and the AILLA recordings give a base for orthography and teaching materials, but no classroom materials were found. Access conditions for the AILLA corpus should be checked, and any use should involve the Tupar\'i community.
\end{document}